\documentclass[12pt,a4paper]{article}
\usepackage[utf8]{inputenc}
\usepackage[russian]{babel}
\usepackage{amsmath}
\usepackage{graphicx}
\usepackage{hyperref}
\usepackage{listings}
\usepackage{geometry}
\geometry{top=2cm,bottom=2cm,left=3cm,right=1.5cm}

\title{Курсовой проект по курсу <<Администрирование ОС Linux>> \\
\large Тема №2: Настройка балансировщика нагрузки HAProxy}
\author{Студенты Кириленко К.С., Лазаренко М.Э., Ивахненко Е.В. 7группы}
\date{2026}

\begin{document}

\maketitle
\newpage

\tableofcontents
\newpage

\section{Введение}
В современном мире веб-сервисы должны обладать высокой доступностью (High Availability) и способностью масштабироваться под возрастающие нагрузки. Одним из ключевых инструментов для достижения этих целей является балансировка нагрузки...

\section{Постановка задачи}
Целью данного курсового проекта является развертывание отказоустойчивого веб-сервиса с использованием балансировщика нагрузки HAProxy и трех бэкенд-серверов Nginx. 
Задачи:
\begin{itemize}
    \item Развернуть изолированную сетевую архитектуру Docker.
    \item Настроить балансировку по алгоритму Round Robin.
    \item Реализовать сохранение сессий (Sticky Sessions).
    \item Обеспечить безопасность и изоляцию контейнеров (Non-root, Read-only rootfs).
    \item Развернуть мониторинг (Prometheus, Grafana, cAdvisor).
    \item Провести автоматизированное Chaos-тестирование.
\end{itemize}

\section{Архитектура системы}
Инфраструктура состоит из трех сетевых контуров:
1. \textbf{frontend-net} -- сеть связи внешнего клиента и HAProxy.
2. \textbf{backend-net} -- изолированная сеть для взаимодействия HAProxy и бэкендов Nginx.
3. \textbf{monitoring-net} -- сеть мониторинга для сбора метрик с помощью Prometheus и визуализации в Grafana.

\section{Техническая реализация и настройка}
\subsection{Конфигурация HAProxy}
HAProxy считывает конфигурацию из файла \texttt{haproxy.conf}. Он принимает запросы на порту \texttt{8080} и перенаправляет их на бэкенды Nginx с использованием Proxy Protocol...

\subsection{Конфигурация Nginx}
Бэкенд-серверы Nginx настроены на прием Proxy Protocol для извлечения реального IP-адреса клиента...

\section{Наблюдаемость и мониторинг}
Для мониторинга утилизации системных ресурсов контейнеров используется связка Prometheus и Grafana...

\section{Тестирование отказоустойчивости}
Программа тестирования имитирует падение до двух серверов бэкенда и измеряет задержку переключения трафика...

\section{Заключение}
В ходе работы над курсовым проектом были освоены навыки администрирования Linux систем, автоматизации оркестрации контейнеров и настройки высокодоступных веб-сервисов...

\end{document}
